% =====================================================================
% paper/sections/zvf_iter34.tex
% NOTE (Deep Review 6c172429): headline numbers in this comment block are
% superseded by the strict-LOO audit (zvf_review_6c17_audit.py): 7/12, p=0.07.
%
% Pillar 2 elevation, iteration 34: cross-pillar ZVF x Phase
% discrimination on the 12-anchor frontier set (iter33).
%
% Question: iter30 proved ZVF is a calibrated leading indicator of
% training failure on the variance-mitigation trajectory set.
% iter33 classified 12 frontier anchors into 4 phase classes
% (plateau, saturation, drift, collapse). This iteration asks
% whether the same ZVF-feature decomposition DISCRIMINATES phase
% classes on the frontier set.
%
% Empirical headline (full tables in
% experiments/results/zvf_iter34_phase_scan.tsv,
% zvf_iter34_predictions.tsv, zvf_iter34_confusion.tsv,
% zvf_iter34_feature_importance.tsv, zvf_iter34_collapse_gap.tsv,
% zvf_iter34_summary.tsv, zvf_iter34_discriminant.tsv; figure
% figures/zvf_iter34.pdf):
%
%   * Leave-one-out 4-class phase classifier on 12 anchors
%     achieves 8/12 = 0.667 accuracy (chance 0.25). Saturation is
%     recovered perfectly (5/5), collapse perfectly (1/1); drift
%     2/3; plateau 0/3 (two plateau anchors are confused with
%     drift and one with collapse).
%
%   * Per-feature permutation importance (200-resample bootstrap):
%     zvf_direction (delta_late_minus_early) is the dominant
%     predictor, with delta_acc = +0.113 (vs +0.058 for
%     zvf_slope, +0.058 for zvf_instability, +0.055 for
%     zvf_discriminator, +0.014 for zvf_level). No shuffled
%     feature beats the base accuracy at the 0.05 level.
%
%   * Collapse-vs-rest gap: zvf_level gap is +0.277 (collapse
%     mean 0.333 vs non-collapse 0.056); zvf_discriminator gap
%     is +0.281 (collapse 0.333 vs 0.052). Nemotron-120B
%     (labelled plateau in iter33) is the only non-collapse
%     trace with zero_frac = 0.55 (above the collapse anchor's
%     0.33) -- the iter33 plateau label for Nemotron is the
%     classifier's only phase misclassification with a high
%     zvf_level.
%
% Source: scripts/zvf_iter34_phasescan.py
% =====================================================================

\section{ZVF Cross-Pillar Phase Discrimination on the 12-Anchor Frontier}
\label{sec:zvf-iter34}

The iter30 analysis reported, \emph{on the simulated variance-mitigation
trajectory set}, that ZVF time-series features carry leading
indicator content. iter33 independently classified 12 frontier
anchors into 4 phase classes. iter34 connects the two: it asks
whether the same ZVF-feature decomposition, evaluated on iter33
trace-level summary statistics, can discriminate phase classes
\emph{on the frontier set itself}. The answer is weak: only the
saturation phase is recovered reliably, and the overall accuracy is not
significantly above a label-permutation null.

\subsection{Setup: ZVF proxies at the trace level}
\label{sec:zvf-iter34-setup}

The iter33 phase score table provides 20 columns of trace-level
summary statistics but no per-step ZVF stream. We map the iter30
ZVF feature battery to four iter33 trace-level proxies:

\begin{itemize}
\item \texttt{zvf\_level} $\equiv$ \texttt{zero\_frac} (fraction
of steps whose mean reward is zero --- a reward-trace statistic, not a
group-level ZVF, which these runs did not log).
\item \texttt{zvf\_instability} $\equiv$ \texttt{peak\_frac}
(fraction of steps at peak reward; high values mean the reward
distribution is bunched at one value).
\item \texttt{zvf\_slope} $\equiv$ \texttt{ols\_slope\_per\_step}
(linear drift in reward that drives ZVF slope).
\item \texttt{zvf\_direction} $\equiv$ \texttt{delta\_late\_minus\_early}
(positive = improving; negative = degrading).
\item \texttt{zvf\_discriminator} $\equiv$
\texttt{zvf\_level} $\times \min\{1,\max(0, 1 - \texttt{zvf\_direction})\}$
(combined level and reward direction; the factor is clipped to $[0,1]$, so
for Qwen3.5-27B, with direction $-0.281$, the score equals its level $0.333$).
\end{itemize}

The classification uses a nearest-centroid classifier on standardized
features with leave-one-out (LOO) cross-validation over the 12 anchors.
In strict LOO, standardization and class centroids are computed on the
training fold only, so a held-out anchor whose class is a singleton cannot
be assigned that class.

\subsection{Finding: only saturation is recovered reliably}
\label{sec:zvf-iter34-finding}

Under strict LOO the classifier recovers the true phase for 7 of 12 anchors
(accuracy $0.583$; majority-class baseline $5/12=0.417$; label-permutation
$p=0.07$ over 2,000 shuffles;
\texttt{experiments/results/zvf\_review\_6c17\_iter34\_strict\_loo\_*.tsv},
\texttt{platform\_modal/scripts/zvf\_review\_6c17\_audit.py}). An earlier
revision reported $8/12$ with the collapse anchor ``correctly classified on
every fold''; that is impossible in strict LOO, because when the only collapse
anchor is held out no collapse centroid exists, and the stored predictions
(\texttt{zvf\_iter34\_predictions.tsv}) do not reproduce under the current
script. Per class:

\begin{itemize}
\item \textbf{Saturation: 5/5.} All five saturation anchors
(Qwen3-8B, gpt-oss-20B, Qwen3-30B-MoE, Qwen3-30B-MoE-Inst,
Qwen3-235B-MoE) are recovered. They share \texttt{zvf\_level} $=0$ except
Qwen3-8B ($0.067$), \texttt{zvf\_instability} $0.33$--$0.52$, and
\texttt{zvf\_direction} $\in [0.00, +0.16]$.
\item \textbf{Drift: 2/3.} Qwen3-32B is predicted as saturation.
\item \textbf{Collapse: 0/1.} The lone collapse anchor (Qwen3.5-27B;
\texttt{zvf\_level} $0.333$, \texttt{zvf\_direction} $-0.281$) is predicted as
drift; a singleton class cannot be recovered in strict LOO.
\item \textbf{Plateau: 0/3.} Qwen3.5-4B and DeepSeek-V3.1 are predicted as
drift; Nemotron-120B is predicted as collapse.
\end{itemize}

\subsection{Feature importance}
\label{sec:zvf-iter34-importance}

The earlier revision ranked features by permutation importance computed with
the non-strict classifier (\texttt{zvf\_direction} first, $\Delta
\text{acc}=+0.113$). For every feature, at least $19.5\%$ of shuffles matched or
beat the base accuracy ($p\ge0.195$), which means \emph{no} feature's
contribution is distinguishable from chance; an earlier sentence that read
this as ``significantly above the permutation null'' inverted the test. We
draw no feature-ranking conclusion.

\subsection{Collapse-vs-rest gap: the collapse signature is
ZVF-level + reward-degrading}
\label{sec:zvf-iter34-collapse}

The collapse-vs-rest gap table
(\texttt{experiments/results/zvf\_iter34\_collapse\_gap.tsv}) compares the
single collapse anchor with the eleven others (descriptive; $n=1$ vs $11$):

\begin{itemize}
\item \texttt{zvf\_level}: collapse mean 0.333 vs non-collapse mean
0.056 (gap $+0.277$). The non-collapse max is 0.55
(Nemotron-120B), and the non-collapse min is 0.00.
\item \texttt{zvf\_discriminator}: collapse 0.333 vs non-collapse
0.052 (gap $+0.281$). The non-collapse max is 0.516
(Nemotron-120B again).
\item \texttt{zvf\_direction}: collapse mean $-0.281$ vs
non-collapse mean $+0.004$ (gap $-0.285$). The non-collapse max
is $+0.156$ (Qwen3-30B-MoE).
\item \texttt{zvf\_slope}: collapse $-0.146$ vs non-collapse
$-0.003$ (gap $-0.143$).
\end{itemize}

The collapse anchor is a \emph{reward-degrading} trace with a non-trivial
fraction of all-zero reward steps, but it is not uniquely separated: on
\texttt{zvf\_level} and \texttt{zvf\_discriminator}, Nemotron-120B scores
higher. The Qwen3.5-27B anchor reaches 75\% peak
reward by step 1 then drops to 38\% by step 3 --- an early
collapse with a single peak and persistent low-reward tail.

A residual concern: Nemotron-120B (labelled plateau in iter33)
has \texttt{zvf\_level} = 0.55, the highest zero-reward fraction
in the 12-anchor set, and \texttt{zvf\_discriminator} = 0.516, the
highest combined score. The classifier predicts it as collapse.
The phase labels themselves disagree across this paper: iter33 (used here)
labels Nemotron-120B \texttt{plateau} and Qwen3.5-27B \texttt{collapse}, whereas the
five-anchor three-phase segmentation of Section~\ref{sec:zvf-scaling} labels
Nemotron-120B \texttt{collapse}. The two labelings come from different
segmentation rules on the same traces; until they are reconciled, neither
can serve as ground truth for a collapse detector.

\subsection{Discriminant table: per-phase feature profile}
\label{sec:zvf-iter34-discriminant}

The per-phase discriminant table
(\texttt{experiments/results/zvf\_iter34\_discriminant.tsv})
gives the mean and SD of each ZVF proxy for each of the 4 phase
classes. Saturation traces are the most uniform (\texttt{zvf\_level}
$=0$ except Qwen3-8B at $0.067$); drift traces share
\texttt{zvf\_level} $=0$ with \texttt{zvf\_direction} in
$\{-0.19, -0.09, -0.08\}$; plateau traces have \texttt{zvf\_direction} in
$\{+0.04, -0.02, +0.06\}$ but \texttt{zvf\_level} ranging from $0$ to $0.55$
(Nemotron-120B).

\subsection{Honest scope}
\label{sec:zvf-iter34-honest}

The 12-anchor frontier set is small for 4-class classification,
the strict-LOO accuracy ($0.583$) is not significantly above a label-permutation
null ($p=0.07$), and the per-class cell counts are 1--5. The
ZVF proxies are summary-statistic derivations, not direct per-step
ZVF measurements; the most informative proxy
(\texttt{zvf\_direction}) is also the simplest
(\texttt{delta\_late\_minus\_early}). The iter34 finding is
therefore best read as \emph{saturation is recognisable from reward-trace
summaries; the other phases are not}, rather than as a phase classifier.